\section{Scriptural Date and Location}
\label{sec:date-location}
{\small Each passage the Mass reads or sings from Scripture, once, in
canonical order; beneath each, its attribution, its setting and the critical
horizon. The Entrance Antiphon and the three orations are composed liturgical
texts and carry no biblical date.\par}

\begin{dossiertable}
Communion Antiphon, first option & Ps 119 (118):4--5 & No place named in the
verses & \chronodate{communion-antiphon-a}{\chronologyannotation{communion-antiphon-a}}\\
\cmidrule(lr){1-4}
\dossierprose{The opening strophe of the alphabetical psalm of the law; the
verses name no author and no place, and no traditional date is held for them.
The date is the modern critical boundary for the whole Psalter, within which
no single psalm can be dated securely.}
\midrule
Responsorial Psalm & Ps 145 (144):2--3, 8--9, 17--18; response v.~18a & No
place named &
\chronodate{responsorial-psalm}{\chronologyannotation{responsorial-psalm}}\\
\cmidrule(lr){1-4}
\dossierprose{Headed in the Douay--Rheims ``Praise, for David himself''
(v.~1): an attribution by title, with no traditional date attached. An
alphabetical hymn to ``God my king''. The date is the Psalter's critical
boundary only.}
\midrule
First Reading & Isa 55:6--9 & Isaias, prophesying ``concerning Juda and
Jerusalem'' (1:1); to his own people &
\chronodate{first-reading}{\chronologyannotation{first-reading}}\\
\cmidrule(lr){1-4}
\dossierprose{Traditional attribution: the prophet Isaias, whose ministry the
\work{Catholic Encyclopedia} (1910) bounds by the range shown; the article
assigns no year to the writing of the book or of chapters 40--66. It treats
chapters 54--55 under the heading ``Second Isaias'' and reports the contention
of most modern non-Catholic scholars that those chapters come from an author
``living towards the close of the Babylonian Captivity'', against which it
sets the Pontifical Biblical Commission's decision of 28~June 1908 and
concludes that the author's speaking ``from the point of view of the
Babylonian Captivity'' is no proof that he lived and wrote then. No date is
given here for the critical position.}
\midrule
Gospel & Mt 20:1--16a & Spoken to the disciples on the way to Jerusalem, after
Peter's question (19:27) and before the third prediction of the Passion
(20:17--19) & \chronodate{gospel}{\chronologyannotation{gospel}}\\
\cmidrule(lr){1-4}
\dossierevent{Narrated event: a parable of Jesus, told within the journey
that ends at Jerusalem (20:17--18).}
\cmidrule(lr){1-4}
\dossierprose{Traditional attribution: St~Matthew the Apostle. The Date column
keeps as disputed alternatives the figures the \work{Catholic Encyclopedia}
reports (1911, 1912): its article on the Gospel opens its discussion of the
date with the words ``Ancient ecclesiastical writers are at variance'' and
says that in its own day ``opinion is rather divided''; the figure of a single
year, printed last, is said in the article on the New Testament of the
original text of Matthew, which its author holds to have been Aramaic. The
article on the Gospel lists further modern figures, which the chronology
record does not carry and which are not printed here.}
\midrule
Communion Antiphon, second option & Jn 10:14 & Spoken to the Pharisees after
the healing of the man born blind (9:40--10:6) &
\chronodate{communion-antiphon-b}{\chronologyannotation{communion-antiphon-b}}\\
\cmidrule(lr){1-4}
\dossierevent{Narrated event: the discourse of the good shepherd.}
\cmidrule(lr){1-4}
\dossierprose{Traditional attribution: St~John the Apostle. The
\work{Catholic Encyclopedia} gives the first range shown for ``the Johannine
writings'' (1912) and the second as ``the general opinion'' for the Gospel,
after conceding that ``we possess no certain historical information'' (1910).
The Missal adapts the verse by adding \latin{dicit Dominus} and the noun
\latin{oves}.}
\midrule
Alleluia verse (its basis) & cf.\ Acts 16:14b, an adaptation & Event:
Philippi, ``the chief city of part of Macedonia, a colony'' (16:12), by the
riverside on the sabbath (16:13) &
\chronodate{gospel-acclamation}{\chronologyannotation{gospel-acclamation}}\\
\cmidrule(lr){1-4}
\dossierevent{Narrated event: the conversion of Lydia during St~Paul's second
missionary journey (Acts 15:36--18:22).}
\cmidrule(lr){1-4}
\dossierprose{Traditional attribution: St~Luke. The two event ranges are those
the \work{Catholic Encyclopedia} gives for the second journey in its articles
on St~Paul (1911) and on biblical chronology (1908); the former places
Philippi in that journey without naming Lydia. The composition date is the one its article on Acts (1907) calls most
probable for the completion of the book, two years after Paul's coming to
Rome. The liturgical verse is a petition formed on Luke's narrative sentence
and is not itself a verse of Scripture.}
\midrule
Second Reading & Phil 1:20c--24, 27a & Paul in bonds (1:7, 13), at Rome; to
``all the saints in Christ Jesus who are at Philippi'' (1:1) &
\chronodate{second-reading}{\chronologyannotation{second-reading}}\\
\cmidrule(lr){1-4}
\dossierprose{Traditional attribution: St~Paul, writing from captivity. The
first figure is the year the \work{Catholic Encyclopedia}'s chronology of
St~Paul gives the captivity letters, within the Roman captivity (1911); the
second is the range its article on the letter gives for Paul ``at Rome''
(1911), where it adds that critics disagree whether the letter falls early or
late in the Roman sojourn and that its author holds for late.}
\end{dossiertable}
